\item \textbf{Problema de repaso.} En una pistola de agua, un pistón impulsa el agua a través de un tubo grande de área $A_1$ hacia un tubo más pequeño de área $A_2$, como se muestra en la figura P14.46. El radio del tubo grande es $1.00 \text{ cm}$ y el del tubo pequeño es $1.00 \text{ mm}$. El tubo menor está $3.00 \text{ cm}$ arriba del tubo mayor. (a) Si la pistola se dispara horizontalmente a una altura de $1.50 \text{ m}$, determine el intervalo de tiempo requerido para que el agua viaje desde la boquilla hasta el suelo. Desprecie la resistencia del aire y suponga presión atmosférica de $1.00 \text{ atm}$. (b) Si el alcance deseado del chorro es $8.00 \text{ m}$, ¿con qué rapidez $v_2$ debe salir el agua por la boquilla? (c) ¿A qué rapidez $v_1$ debe moverse el émbolo para lograr el alcance deseado? (d) ¿Cuál es la presión en la boquilla? (e) Encuentre la presión requerida en el tubo mayor. (f) Calcule la fuerza que se debe ejercer sobre el gatillo para lograr el alcance deseado. (La fuerza que se debe ejercer está arriba de la presión atmosférica.)
